\section*{第 \ref{ch:g8:equations} 章——字母运算与方程}

\begin{solution}{exo:g8:equations:1}
$5(2x + 3) = 10x + 15$.

$-3(4x - 1) = -12x + 3$.

$2(3x - 2) + 4(x + 5) = 6x - 4 + 4x + 20 = 10x + 16$.
\end{solution}

\begin{solution}{exo:g8:equations:2}
$(x+2)(x+7) = x^2 + 7x + 2x + 14 = x^2 + 9x + 14$.

$(3x+1)(2x-5) = 6x^2 - 15x + 2x - 5 = 6x^2 - 13x - 5$.

$(x-4)(x-6) = x^2 - 6x - 4x + 24 = x^2 - 10x + 24$.
\end{solution}

\begin{solution}{exo:g8:equations:3}
$4 \times 3 - 5 = 7$：成立。 \quad
$2 \times 3 + 1 = 7$ 与 $3 + 4 = 7$ 均成立。 \quad
$3^2 = 9$ 与 $6 \times 3 - 9 = 9$ 均成立——$x = 3$ 能同时
满足这三个式子。
\end{solution}

\begin{solution}{exo:g8:equations:4}
$x + 9 = 4$：$x = -5$（检验：$-5 + 9 = 4$）。

$5x = -35$: $x = -7$.

$\frac x3 = 8$: $x = 24$.

$2x - 7 = 13$：$2x = 20$，$x = 10$（检验：$20 - 7 = 13$）。
\end{solution}

\begin{solution}{exo:g8:equations:5}
$6x + 5 = 2x + 25$：$4x = 20$，故 $x = 5$。

$9 - 4x = x - 11$：$20 = 5x$，故 $x = 4$。

$3(x + 2) = 2x + 11$：$3x + 6 = 2x + 11$，故 $x = 5$。
\end{solution}

\begin{solution}{exo:g8:equations:6}
$8x - 12 = 5x + 30$，故 $3x = 42$ 且 $x = 14$。检验：
$4(28 - 3) = 100$ 与 $5(14 + 6) = 100$。
\end{solution}

\begin{solution}{exo:g8:equations:7}
“$3x - 8 = 19$”：$3x = 27$，$x = 9$。

“$2x + 5 = x + 12$”：$x = 7$。
\end{solution}

\begin{solution}{exo:g8:equations:8}
宽为 $x$，长为 $x + 5$；\omterm{def:g6:measure:perimeter}{周长} 
$2\bigl(x + (x + 5)\bigr) = 4x + 10 = 58$，故 $4x = 48$ 且 $x = 12$：
宽为 $12$ cm，长为 $17$ cm（检验：
$2 \times (12 + 17) = 58$）。
\end{solution}

\begin{solution}{exo:g8:equations:9}
这三个数分别为 $n - 1$、$n$、$n + 1$，其和为 $3n = 141$，故
$n = 47$：这三个数是 $46$、$47$、$48$。
\end{solution}

\begin{solution}{exo:g8:equations:10}
$3x + 4 > 19$：$3x > 15$，故 $x > 5$。

$7 - 2x \geq 1$：$-2x \geq -6$，两边除以 $-2$，不等号方向改变：
$x \leq 3$.

$5(x-1) < 2x + 7$：$5x - 5 < 2x + 7$，故 $3x < 12$ 且 $x < 4$。
\end{solution}

\begin{solution}{exo:g8:equations:11}
经过 $n$ 个月后，米娅有 $140 + 15n$，汤姆有 $200 + 10n$。米娅领先
的条件为
\[
140 + 15n > 200 + 10n
\ \Longleftrightarrow\
5n > 60
\ \Longleftrightarrow\
n > 12 :
\]
即从第 $13$ 个月起。
\end{solution}

\begin{solution}{exo:g8:equations:12}
两边同乘 $20$：
\[
5(x + 3) = 4(2x - 1)
\ \Longrightarrow\
5x + 15 = 8x - 4
\ \Longrightarrow\
19 = 3x
\ \Longrightarrow\
x = \frac{19}{3}.
\]
检验：$\frac{19/3 + 3}{4} = \frac{28/3}{4} = \frac{28}{12} =
\frac73$，而 $\frac{2 \times 19/3 - 1}{5} = \frac{35/3}{5} =
\frac{35}{15} = \frac73$。两者相等，方程的解为 $\frac{19}{3}$。
\end{solution}

\begin{solution}{pb:g8:equations:1}
\textbf{1.} 两次运用分配律
（由\cref{thm:g8:equations:expand}），取 $c = a$，$d = b$：
\[
(a + b)(a + b) = a^2 + ab + ba + b^2 = a^2 + 2ab + b^2 ,
\]
因为 $ab$ 与 $ba$ 是同一个积。

\textbf{2.} 展开方法相同，注意符号：
\[
(a - b)(a - b) = a^2 - ab - ba + b^2 = a^2 - 2ab + b^2 ,
\]
\[
(a + b)(a - b) = a^2 - ab + ba - b^2 = a^2 - b^2 :
\]
这一次交叉项 $-ab$ 与 $+ba$ 相互抵消，
而不是加倍。

\textbf{3.} 边长为 $a + b$ 的正方形分成四块：一个
边长为 $a \times a$ 的正方形（即 $a^2$），一个边长为 $b \times b$ 的
正方形（即 $b^2$），以及\emph{两个}长方形 $a \times b$（合起来，即
$2ab$）。恒等式 $(a+b)^2 = a^2 + 2ab + b^2$ 表明
四块的\omterm{def:g6:measure:area}{面积}恰好铺满大正方形——这正是
两边以同样方式拆分时双重分配律的图示。

\textbf{4.} 利用上述恒等式：
\begin{align*}
101^2 &= (100 + 1)^2 = 10\,000 + 2 \times 100 + 1 = 10\,201, \\
99^2 &= (100 - 1)^2 = 10\,000 - 200 + 1 = 9\,801, \\
49 \times 51 &= (50 - 1)(50 + 1) = 50^2 - 1^2 = 2\,499 .
\end{align*}

\textbf{5.} 取 $a = 3$，$b = 4$：$(3 + 4)^2 = 49$，而
$3^2 + 4^2 = 9 + 16 = 25$。佐伊的公式漏掉了图中的两个
$a \times b$ 长方形——即缺少的
$2ab = 24$，且 $25 + 24 = 49$。

\textbf{6.} $1$；$1 + 3 = 4$；$1 + 3 + 5 = 9$；
$1 + 3 + 5 + 7 = 16$；$1 + 3 + 5 + 7 + 9 = 25$。结果是
完全平方数 $1^2, 2^2, 3^2, 4^2, 5^2$。

\textbf{7.} \omterm{def:g3:numbers:evenodd}{奇数}从 $1$ 开始，每次增加
$2$：第 $n$ 个\omterm{def:g3:numbers:evenodd}{奇数}经过 $n - 1$ 步、每步增加 $2$，从
$1$ 出发，故它等于 $1 + 2(n - 1) = 2n - 1$。检验：$n = 1, 2, 3$
时分别得 $1, 3, 5$。

\textbf{8.} 由第一个恒等式，
\[
(n + 1)^2 - n^2 = n^2 + 2n + 1 - n^2 = 2n + 1 :
\]
相邻两个平方数之差，恰好是
第 $(n+1)$ 个\omterm{def:g3:numbers:evenodd}{奇数}（第 7 题中以 $n + 1$ 代替 $n$：
$2(n+1) - 1 = 2n + 1$）。

\textbf{9.} 依次拼出各正方形。从一个方格开始，
即 $1 = 1^2$。由边长为 $1$ 的正方形过渡到
边长为 $2$ 的正方形，需添加 $2 \times 1 + 1 = 3$ 个方格（第
8 题）：$1 + 3 = 2^2$。过渡到边长为 $3$ 的正方形，需添加 $2 \times 2 + 1 = 5$
个方格：$1 + 3 + 5 = 3^2$。每个新增的\omterm{def:g3:numbers:evenodd}{奇数}恰好是
补成下一个正方形的 L 形外层，因此
添加前 $n$ 个\omterm{def:g3:numbers:evenodd}{奇数}后，边长为 $n$ 的正方形便拼成：
$1 + 3 + \dots + (2n - 1) = n^2$。

\textbf{10.} 从 $1$ 到 $99$ 的\omterm{def:g3:numbers:evenodd}{奇数}恰为其中的前 $50$ 个
（由 $2n - 1 = 99$ 得 $n = 50$），故
\[
1 + 3 + 5 + \dots + 99 = 50^2 = 2\,500 .
\]
又 $1 + 3 + \dots + (2n - 1) = 400 = 20^2$ 要求 $n = 20$：
即前\emph{二十}个\omterm{def:g3:numbers:evenodd}{奇数}（从 $1$ 到 $39$）之和为
$400$。

\textbf{11.} $x^2 + 6x + 9 = x^2 + 2 \times 3 \times x + 3^2 =
(x + 3)^2$：由第一个恒等式，取 $a = x$，$b = 3$。于是
方程化为 $(x + 3)^2 = 0$。一个平方为零，仅当
数本身为零（由\cref{thm:g8:negprod:rules}：非零数的距离不为零，
故其平方的距离亦不为零）。
于是 $x + 3 = 0$，唯一解为 $x = -3$。

\textbf{12.} 由\cref{thm:g8:negprod:rules}，乘积的距离
等于各因式距离之积。若两个因式都不为零，则它们的距离都不为零，
于是乘积的距离不为零，
即乘积不可能为 $0$。因此零乘积至少有一个因式为零。
对于 $(x - 2)(x + 5) = 0$，必有
$x - 2 = 0$ 或 $x + 5 = 0$，得两个解 $x = 2$ 与
$x = -5$。

\textbf{13.} $x^2 = 49$ 化为 $x^2 - 49 = 0$，再用第三个
恒等式将其因式分解：
\[
x^2 - 49 = x^2 - 7^2 = (x - 7)(x + 7) = 0 .
\]
由第 12 题，$x = 7$ 或 $x = -7$：\emph{两个}解。本章此前所解的每个
方程都是一次方程，且
恰有一个解；方程中的平方项可以
产生两个解。

\textbf{14.} 将所有项移到左边：
$x^2 - 6x + 9 = 0$，并识别出第二个恒等式：
$x^2 - 2 \times 3 \times x + 3^2 = (x - 3)^2$。于是方程为
$(x - 3)^2 = 0$，由此得 $x - 3 = 0$（同第 11 题）：
$x = 3$ 是唯一解——\cref{exo:g8:equations:3}中的检验
已找到它，而因式分解证明
不存在其他解。

\textbf{15.} 由第三个恒等式，
\[
2026^2 - 2025^2 = (2026 - 2025)(2026 + 2025)
= 1 \times 4051 = 4\,051 .
\]
一般地，$(n+1)^2 - n^2 = (n + 1 - n)(n + 1 + n) = 2n + 1$：
相邻两个平方数之差恰为这两数之和——
与第 8 题相同的 $2n + 1$，此次由
第三个恒等式得到，而非第一个。
\end{solution}
